%--------------------------------------------------------------------------------------
% Page layout setup
%--------------------------------------------------------------------------------------
% we need to redefine the pagestyle plain
% another possibility is to use the body of this command without \fancypagestyle
% and use \pagestyle{fancy} but in that case the special pages
% (like the ToC, the References, and the Chapter pages)remain in plane style

\pagestyle{plain}

\setcounter{tocdepth}{3}
%\sectionfont{\large\upshape\bfseries}
\setcounter{secnumdepth}{3}

\sloppy % Margón túllógó sorok tiltása.
\widowpenalty=10000 \clubpenalty=10000 %A fattyú- és árvasorok elkerülése
\def\hyph{-\penalty0\hskip0pt\relax} % Kötőjeles szavak elválasztásának engedélyezése


%--------------------------------------------------------------------------------------
% Setup hyperref package
%--------------------------------------------------------------------------------------
\hypersetup{
    % bookmarks=true,            % show bookmarks bar?
    unicode=true,              % non-Latin characters in Acrobat's bookmarks
    pdftitle={\vikcim},        % title
    pdfauthor={\szerzoMeta},    % author
    pdfsubject={\vikdoktipus}, % subject of the document
    pdfcreator={\szerzoMeta},   % creator of the document
    pdfproducer={},    % producer of the document
    pdfkeywords={},    % list of keywords (separate then by comma)
    pdfnewwindow=true,         % links in new window
    colorlinks=true,           % false: boxed links; true: colored links
    linkcolor=black,           % color of internal links
    citecolor=black,           % color of links to bibliography
    filecolor=black,           % color of file links
    urlcolor=black             % color of external links
}


%--------------------------------------------------------------------------------------
% Set up listings
%--------------------------------------------------------------------------------------
\definecolor{lightgray}{rgb}{0.95,0.95,0.95}
\lstset{
	basicstyle=\scriptsize\ttfamily, % print whole listing small
	keywordstyle=\color{black}\bfseries, % bold black keywords
	identifierstyle=, % nothing happens
	% default behavior: comments in italic, to change use
	% commentstyle=\color{green}, % for e.g. green comments
	stringstyle=\scriptsize,
	showstringspaces=false, % no special string spaces
	aboveskip=3pt,
	belowskip=3pt,
	backgroundcolor=\color{lightgray},
	columns=flexible,
	keepspaces=true,
	escapeinside={(*@}{@*)},
	captionpos=b,
	breaklines=true,
	frame=single,
	float=!ht,
	tabsize=2,
	literate=*
		{á}{{\'a}}1	{é}{{\'e}}1	{í}{{\'i}}1	{ó}{{\'o}}1	{ö}{{\"o}}1	{ő}{{\H{o}}}1	{ú}{{\'u}}1	{ü}{{\"u}}1	{ű}{{\H{u}}}1
		{Á}{{\'A}}1	{É}{{\'E}}1	{Í}{{\'I}}1	{Ó}{{\'O}}1	{Ö}{{\"O}}1	{Ő}{{\H{O}}}1	{Ú}{{\'U}}1	{Ü}{{\"U}}1	{Ű}{{\H{U}}}1
}


%--------------------------------------------------------------------------------------
% Set up theorem-like environments
%--------------------------------------------------------------------------------------
% Using ntheorem package -- see http://www.math.washington.edu/tex-archive/macros/latex/contrib/ntheorem/ntheorem.pdf

\theoremstyle{plain}
\theoremseparator{.}
\newtheorem{example}{\pelda}

\theoremseparator{.}
%\theoremprework{\bigskip\hrule\medskip}
%\theorempostwork{\hrule\bigskip}
\theorembodyfont{\upshape}
\theoremsymbol{{\large \ensuremath{\centerdot}}}
\newtheorem{definition}{\definicio}

\theoremseparator{.}
%\theoremprework{\bigskip\hrule\medskip}
%\theorempostwork{\hrule\bigskip}
\newtheorem{theorem}{\tetel}


%--------------------------------------------------------------------------------------
% Some new commands and declarations
%--------------------------------------------------------------------------------------
\newcommand{\code}[1]{{\upshape\ttfamily\scriptsize\indent #1}}
\newcommand{\doi}[1]{DOI: \href{http://dx.doi.org/\detokenize{#1}}{\raggedright{\texttt{\detokenize{#1}}}}} % A hivatkozások közt így könnyebb DOI-t megadni.

\DeclareMathOperator*{\argmax}{arg\,max}
%\DeclareMathOperator*[1]{\floor}{arg\,max}
\DeclareMathOperator{\sign}{sgn}
\DeclareMathOperator{\rot}{rot}


%--------------------------------------------------------------------------------------
% Setup captions
%--------------------------------------------------------------------------------------
\captionsetup[figure]{
	width=.75\textwidth,
	aboveskip=10pt}

\renewcommand{\captionlabelfont}{\bf}
%\renewcommand{\captionfont}{\footnotesize\it}

%--------------------------------------------------------------------------------------
% Hyphenation exceptions
%--------------------------------------------------------------------------------------
\hyphenation{Shakes-peare Mar-seilles ár-víz-tű-rő tü-kör-fú-ró-gép}


\author{\vikszerzo}
\title{\viktitle}

%--------------------------------------------------------------------------------------
% Todo notes
%--------------------------------------------------------------------------------------
% The page margins are narrow (see \marginsize in thesis.tex), so margin notes get
% squeezed. Prefer \todo[inline]{...} for anything longer than a few words.
\setlength{\marginparwidth}{13mm}
\newcommand{\todoinline}[1]{\todo[inline]{#1}}


%--------------------------------------------------------------------------------------
% Acronyms (acro)
%--------------------------------------------------------------------------------------
\acsetup{
  make-links = true,     % \ac links to the list of acronyms
  list/heading = chapter*,
}
% Declare every acronym here; use \ac{key} in the text (first use expands it),
% \acs{key} for the short form, \acl{key} for the long form, \acp{key} for plural.
% Print the list with \printacronyms (see thesis.tex).
\DeclareAcronym{vc}{
  short = VC,
  long  = verifiable credential,
}
\DeclareAcronym{zkp}{
  short = ZKP,
  long  = zero-knowledge proof,
}



%--------------------------------------------------------------------------------------
% Colored boxes (tcolorbox)
%--------------------------------------------------------------------------------------
% theo: top-ruled box with a detached title, from Bertalan's BPM paper.
% Usage: \begin{tcolorbox}[theo=blue, title={Key insight}] ... \end{tcolorbox}
\tcbset{%
    theo/.style={%
        enhanced,
        breakable = true,
        sharp corners,
        before skip=20pt,
        toprule=1pt, rightrule=0pt, bottomrule=0pt, leftrule=0pt,
        colback=#1!15, colframe=#1!50, coltitle=#1!60!black,
        detach title,
        overlay unbroken and first ={
            \node[minimum width=1cm, anchor=south west, font=\bfseries]
               at (frame.north west) {\tcbtitle};
        }
    }
}

% Ready-made boxes for the common cases. The optional argument is the title.
\newtcolorbox{infobox}[1][]{enhanced, breakable, colback=blue!5, colframe=blue!50!black, fonttitle=\bfseries, title={#1}}
\newtcolorbox{warnbox}[1][]{enhanced, breakable, colback=orange!8, colframe=orange!70!black, fonttitle=\bfseries, title={#1}}
\newtcolorbox{keybox}[1][]{theo=teal, title={#1}}


%--------------------------------------------------------------------------------------
% Cross-references (cleveref)
%--------------------------------------------------------------------------------------
% Write \cref{fig:overview} instead of Figure~\ref{fig:overview}; \Cref at sentence start.
% Names for the ntheorem environments defined above; cleveref knows the standard ones.
\crefname{definition}{Definition}{Definitions}
\crefname{theorem}{Theorem}{Theorems}
\crefname{example}{Example}{Examples}
\crefname{lstlisting}{Listing}{Listings}
\crefname{rq}{RQ}{RQs}

% Research questions: \begin{rqlist} \item\label{rq:scale} ... \end{rqlist}, then \cref{rq:scale}.
\newlist{rqlist}{enumerate}{1}
\setlist[rqlist]{label=\textbf{RQ\arabic*}, ref=\arabic*, leftmargin=*}
\crefalias{rqlisti}{rq}
\creflabelformat{rq}{#2#1#3}
